%Data institusi

%\input{Dozentdaten}

%\renewcommand{\fachbereich}{Fachbereich}

%\renewcommand{\dozent}{Prof. Dr. . }

%Data ujian

\renewcommand{\klausurgebiet}{ }

\renewcommand{\klausurtyp}{ }

\renewcommand{\klausurdatum}{ . 20}

\klausurvorspann {\fachbereich} {\klausurdatum} {\dozent} {\klausurgebiet} {\klausurtyp}

%Data untuk tabel nilai berikut

\renewcommand{\aeins}{  3  }

\renewcommand{\azwei}{  3   }

\renewcommand{\adrei}{  0  }

\renewcommand{\avier}{  4  }

\renewcommand{\afuenf}{  5  }

\renewcommand{\asechs}{  4  }

\renewcommand{\asieben}{  9  }

\renewcommand{\aacht}{  0  }

\renewcommand{\aneun}{  8  }

\renewcommand{\azehn}{  6  }

\renewcommand{\aelf}{  0  }

\renewcommand{\azwoelf}{  3  }

\renewcommand{\adreizehn}{  3   }

\renewcommand{\avierzehn}{  0  }

\renewcommand{\afuenfzehn}{  0  }

\renewcommand{\asechzehn}{  7  }

\renewcommand{\asiebzehn}{  55  }

\renewcommand{\aachtzehn}{    }

\renewcommand{\aneunzehn}{    }

\renewcommand{\azwanzig}{    }

\renewcommand{\aeinundzwanzig}{    }

\renewcommand{\azweiundzwanzig}{    }

\renewcommand{\adreiundzwanzig}{    }

\renewcommand{\avierundzwanzig}{    }

\renewcommand{\afuenfundzwanzig}{    }

\renewcommand{\asechsundzwanzig}{    }

\punktetabellesechzehn

\klausurnote

\newpage

\setcounter{section}{K}

__NOEDITSECTION__

\inputaufgabeklausurloesung
{3}

{

Definisikan istilah-istilah berikut
\zusatzklammer {yang dicetak miring} {} {.}
\aufzaehlungsechs{
\stichwort {Kelengkungan Gauss} {}
dari suatu permukaan diferensiabel

\mavergleichskette
{\vergleichskette
{ Y
}
{ \subseteq }{ \R^3
}
{  }{
}
{    }{
}
{   }{
}
}
{}{}{}
di suatu titik

\mavergleichskette
{\vergleichskette
{ P
}
{ \in }{ Y
}
{  }{
}
{    }{
}
{   }{
}
}
{}{}{.}

}{Suatu
\stichwort {manifold topologis} {}
berdimensi $n$.

}{Suatu
\stichwort {medan vektor tak bergantung waktu} {}
pada manifold diferensiabel $M$.

}{Suatu
\stichwort {bentuk diferensial berderajat $k$} {}
pada manifold diferensiabel $M$.

}{Suatu \stichwort {manifold Riemann} {.}

}{Suatu koneksi linear yang
\stichwort {bebas torsi} {}
pada bundel tangen manifold Riemann $M$.
}

}
{

\aufzaehlungsechs{Hasil kali
\mathl{\kappa_1 \cdot \kappa_2}{} dari kedua
\definitionsverweis {kelengkungan utama}{}{}
di $P$ disebut
\definitionswort {kelengkungan Gauss}{}
permukaan $Y$ di $P$.
}{Suatu
\definitionsverweis {ruang Hausdorff}{}{}
\definitionsverweis {topologis}{}{}
$M$ disebut manifold topologis berdimensi $n$ jika terdapat suatu
\definitionsverweis {penutup terbuka}{}{}

\mavergleichskette
{\vergleichskette
{ M
}
{ = }{  \bigcup_{i \in I} U_i
}
{  }{
}
{    }{
}
{   }{
}
}
{}{}{}
sedemikian sehingga setiap $U_i$
\definitionsverweis {homeomorfik}{}{}
dengan suatu
\definitionsverweis {subhimpunan terbuka}{}{}
dari $\R^n$.
}{Suatu pemetaan
\maabbdisp {F} {M} {TM
} {}
dengan sifat

\mavergleichskette
{\vergleichskette
{ F(P)
}
{ \in }{  T_PM
}
{  }{
}
{    }{
}
{   }{
}
}
{}{}{}
untuk setiap titik

\mavergleichskette
{\vergleichskette
{ P
}
{ \in }{ M
}
{  }{
}
{    }{
}
{   }{
}
}
{}{}{}
disebut medan vektor
\zusatzklammer {tak bergantung waktu} {} {}.
}{Suatu
$k$-bentuk diferensial
adalah suatu
\definitionsverweis {penampang}{}{}
dalam
\definitionsverweis {hasil kali baji}{}{}
ke-$k$ dari
\definitionsverweis {bundel kotangen}{}{.}
}{Suatu
\definitionsverweis {manifold diferensiabel}{}{}
$M$ disebut manifold Riemann jika pada setiap
\definitionsverweis {ruang tangen}{}{}

\mathbed {T_PM} {}
 {P \in M} {}
 {} {} {} {,}
diberikan suatu
\definitionsverweis {hasil kali skalar}{}{}

\mathl{\left\langle - , -  \right\rangle_P}{} sedemikian sehingga, untuk setiap peta
\maabbdisp {\alpha} { U } { V
} {}
dengan
\mathl{V\subseteq \R^n}{}, fungsi-fungsi
\zusatzklammer {untuk \mathlk{1 \leq i,j \leq n}{}} {} {}
\maabbeledisp {g_{ij}} { V } {\R
} { Q } {g_{ij}(Q) =  \left\langle  T(\alpha^{-1})(e_i)  ,  T(\alpha^{-1}) (e_j)   \right\rangle_{ \alpha^{-1}(Q) }
} {,}
bersifat $C^1$-\definitionsverweis {terdiferensialkan}{}{}.
}{Koneksi tersebut disebut
bebas torsi
jika

\mavergleichskettedisp
{\vergleichskette
{ \nabla_V W- \nabla_W V
}
{ =} { [V,W]
}
{ } {
}
{ } {
}
{ } {
}
}
{}{}{}
berlaku untuk sebarang medan vektor $V,W$.
}

 }

__NOEDITSECTION__

\inputaufgabeklausurloesung
{3}

{

Nyatakan teorema-teorema berikut.
\aufzaehlungdrei{
\stichwort {Teorema eksistensi transpor paralel pada hiperpermukaan} {.}}{Teorema tentang deskripsi lokal bentuk diferensial.}{
\stichwort {Versi balok} {}
dari teorema Stokes.}

}
{

\aufzaehlungdrei{\faktsituation {Misalkan

\mavergleichskette
{\vergleichskette
{ Y
}
{ \subseteq }{ W
}
{ \subseteq }{\R^n
}
{    }{
}
{   }{
}
}
{}{}{,}
$W$
\definitionsverweis {terbuka}{}{,}
suatu
\definitionsverweis {hiperpermukaan diferensiabel}{}{}
dan misalkan
\maabb {\gamma} { I } { Y
} {}
suatu
\definitionsverweis {kurva diferensiabel}{}{.}
Misalkan

\mavergleichskette
{\vergleichskette
{ t_0
}
{ \in }{ I
}
{  }{
}
{    }{
}
{   }{
}
}
{}{}{,}

\mavergleichskette
{\vergleichskette
{ \gamma(t_0)
}
{ = }{ P
}
{  }{
}
{    }{
}
{   }{
}
}
{}{}{}
dan misalkan

\mavergleichskette
{\vergleichskette
{ v
}
{ \in }{ T_PY
}
{  }{
}
{    }{
}
{   }{
}
}
{}{}{}
suatu
\definitionsverweis {vektor tangen}{}{.}}
\faktfolgerung {Maka terdapat suatu
\definitionsverweis {medan vektor tangensial}{}{}
$F$ sepanjang $\gamma$ yang ditentukan secara unik,
\definitionsverweis {paralel}{}{}
dan memenuhi

\mavergleichskettedisp
{\vergleichskette
{ F(t_0)
}
{ =} { v
}
{ } {
}
{ } {
}
{ } {
}
}
{}{}{.}}
\faktzusatz {}
\faktzusatz {}}{Misalkan $M$ suatu
\definitionsverweis {manifold diferensiabel}{}{}
dan
\mathl{U \subseteq M}{} suatu
\definitionsverweis {subhimpunan terbuka}{}{}
dengan peta
\maabbdisp {\alpha} { U } { V
} {}
dan
\mathl{V \subseteq \R^n}{} terbuka. Misalkan
\maabbdisp {x_j} { U } {\R
} {} adalah fungsi-fungsi koordinat terkait,
\mathl{1 \leq j \leq n}{.} Maka setiap bentuk pada $U$ yang merupakan $k$-\definitionsverweis {bentuk diferensial}{}{}

\mathl{\omega \in
{ \mathcal E }^{ k } (  U   )}{} dapat ditulis secara unik sebagai

\mavergleichskettedisp
{\vergleichskette
{  \omega
}
{ =} { \sum_{J,\,  { \# \left(   J  \right) } = k}   f_J dx_J
}
{ } {
}
{ } {
}
{ } {
}
}
{}{}{}
dengan fungsi-fungsi yang ditentukan secara unik
\maabbdisp {f_J} { U } {\R
} {.}}{\faktsituation {Misalkan

\mavergleichskette
{\vergleichskette
{ Q
}
{ \subseteq }{ \R^n
}
{  }{
}
{    }{
}
{   }{
}
}
{}{}{}
suatu balok $n$-dimensional yang sejajar sumbu
\zusatzklammer {dengan sisi tetapi tanpa rusuk} {} {}
dengan
\definitionsverweis {batas}{}{}

\mathl{\partial Q}{} dan $\omega$ suatu
\definitionsverweis {bentuk diferensial}{}{}
yang didefinisikan pada $Q$, berderajat $(n-1)$, dan
\definitionsverweis {terdiferensialkan secara kontinu}{}{}
.}
\faktfolgerung {Maka

\mavergleichskettedisp
{\vergleichskette
{
\int_{  Q  }  d \omega
}
{ =} {
\int_{  \partial Q  }   \omega
}
{ } {
}
{ } {
}
{ } {
}
}
{}{}{.}}
\faktzusatz {}
\faktzusatz {}}

 }

__NOEDITSECTION__

\inputaufgabeklausurloesung
{0}

{

}
{  }

__NOEDITSECTION__

\inputaufgabeklausurloesung
{4 (1+3)}

{

Misalkan $M$ elips yang diberikan oleh

\mavergleichskettedisp
{\vergleichskette
{ x^2+3y^2
}
{ =} { 2
}
{ } {
}
{ } {
}
{ } {
}
}
{}{}{}
dalam $\R^2$,

\mavergleichskette
{\vergleichskette
{ P
}
{ = }{ \begin{pmatrix}  1  \\ { \frac{   1  }{  \sqrt{3}  } }    \end{pmatrix}
}
{  }{
}
{    }{
}
{   }{
}
}
{}{}{}
dan

\mavergleichskette
{\vergleichskette
{ Q
}
{ = }{ \begin{pmatrix}  -1 \\ { \frac{   1  }{  \sqrt{3}  } }    \end{pmatrix}
}
{  }{
}
{    }{
}
{   }{
}
}
{}{}{}
adalah titik-titik pada $M$.
\aufzaehlungdrei{Tunjukkan bahwa

\mavergleichskettedisp
{\vergleichskette
{ v
}
{ =} { \begin{pmatrix}  -  \sqrt{3}  \\ 1    \end{pmatrix}
}
{ } {
}
{ } {
}
{ } {
}
}
{}{}{}
merupakan
\definitionsverweis {vektor tangen}{}{}
di $P$ pada $M$.
}{Tentukan citra $v$ di bawah suatu transpor paralel pada $M$ dari $P$ ke $Q$.
}{
}

}
{

\aufzaehlungzwei {Ruang tangen di titik $\begin{pmatrix}  x  \\ y   \end{pmatrix}$ tegak lurus terhadap gradien
\mathl{\begin{pmatrix}  2x \\ 6y   \end{pmatrix}}{,} sehingga di $P$ ruang tersebut tegak lurus terhadap

\mavergleichskettedisp
{\vergleichskette
{ \begin{pmatrix}  2 \\ { \frac{   6 }{  \sqrt{3}  } }    \end{pmatrix}
}
{ =} { 2 \begin{pmatrix}  1  \\ \sqrt{3}    \end{pmatrix}
}
{ } {
}
{ } {
}
{ } {
}
}
{}{}{.}
Vektor $v$ memenuhi syarat ini.
} {Transpor paralel sepanjang suatu lintasan $\gamma$ pada $M$ dari $P$ ke $Q$ merupakan isometri
\maabb {} {T_PM} {T_QM
} {.}
Ruang tangen di $Q$ tegak lurus terhadap
\mathl{\begin{pmatrix}  - 1  \\ \sqrt{3}    \end{pmatrix}}{} dan direntang oleh
\mathl{\begin{pmatrix}  \sqrt{3}  \\ 1    \end{pmatrix}}{.} Karena normanya tidak berubah, hasilnya hanya mungkin
\mathl{\pm \begin{pmatrix}  \sqrt{3}  \\ 1    \end{pmatrix}}{}. Karena $v$ menunjuk berlawanan arah jarum jam pada $M$, hasilnya juga harus demikian
\zusatzklammer {misalnya berdasarkan
[[Hyperfläche/Geodätische Kurve/Tangentiale Richtung/Paralleltransport/Aufgabe|Soal 6.9 (Geometri diferensial (Osnabrück 2023))]]} {} {,}
sehingga hasilnya adalah
\mathl{- \begin{pmatrix}  \sqrt{3}  \\ 1    \end{pmatrix}}{}.
}

 }

__NOEDITSECTION__

\inputaufgabeklausurloesung
{5}

{

Buktikan teorema tentang orientasi-orientasi pada hiperpermukaan diferensiabel.

}
{

Menurut
[[Hyperfläche/Glatt/Normalenfeld/Fakt|Lema 2.4 (Geometri diferensial (Osnabrück 2023))]],
terdapat suatu medan normal satuan yang merepresentasikan sebuah orientasi, dan negatifnya menghasilkan orientasi lain. Kita namakan keduanya $G$ dan $-G$. Sekarang, misalkan
\maabbdisp {N} { Y } { \R^n
} {}
suatu medan normal satuan kontinu sebarang. Kita tinjau himpunan-himpunan tempat keduanya berimpit

\mavergleichskettedisp
{\vergleichskette
{ Y_1
}
{ =} { \left\{ P \in Y  \mid   N(P) = G(P)         \right\}
}
{ } {
}
{ } {
}
{ } {
}
}
{}{}{}
dan

\mavergleichskettedisp
{\vergleichskette
{ Y_2
}
{ =} { \left\{ P \in Y  \mid   N(P) = - G(P)         \right\}
}
{ } {
}
{ } {
}
{ } {
}
}
{}{}{.}
Karena untuk setiap titik

\mavergleichskette
{\vergleichskette
{ P
}
{ \in }{ Y
}
{  }{
}
{    }{
}
{   }{
}
}
{}{}{}
hanya terdapat dua nilai yang mungkin,

\mavergleichskettedisp
{\vergleichskette
{ N(P)
}
{ =} { \pm G(P)
}
{ } {
}
{ } {
}
{ } {
}
}
{}{}{}
maka $Y$ adalah gabungan saling lepas dari
\mathkor {} {Y_1} {dan} {Y_2} {.}
Sebagai himpunan tempat fungsi-fungsi kontinu berimpit,
\mathkor {} {Y_1} {dan} {Y_2} {}
\definitionsverweis {tertutup}{}{}
\zusatzklammer {dan dengan demikian juga
\definitionsverweis {terbuka}{}{}
dalam $Y$} {} {.}
Karena keterhubungan, diperoleh

\mavergleichskette
{\vergleichskette
{ Y
}
{ = }{ Y_1
}
{  }{
}
{    }{
}
{   }{
}
}
{}{}{}
atau

\mavergleichskette
{\vergleichskette
{ Y
}
{ = }{ Y_2
}
{  }{
}
{    }{
}
{   }{
}
}
{}{}{.}

 }

__NOEDITSECTION__

\inputaufgabeklausurloesung
{4}

{

Misalkan

\mavergleichskette
{\vergleichskette
{ Y
}
{ \subseteq }{ W
}
{ \subseteq }{ \R^3
}
{    }{
}
{   }{
}
}
{}{}{,}
$W$
\definitionsverweis {terbuka}{}{,}
dan suatu
\definitionsverweis {permukaan diferensiabel}{}{}
yang dilengkapi dengan
\definitionsverweis {medan normal satuan}{}{}
$N$. Misalkan
\maabb {\gamma} { I } { Y
} {}
suatu kurva terdiferensialkan secara kontinu dan
\maabbdisp {F} { I } { \R^3
} {}
suatu
\definitionsverweis {medan vektor tangensial}{}{}
diferensiabel sepanjang $\gamma$ yang
\definitionsverweis {paralel}{}{.}
Tunjukkan bahwa medan vektor
\maabbeledisp {} { I } { \R^3
} { t } { N(\gamma(t)) \times F(t)
} {,}
juga paralel, dengan $\times$ menyatakan
\definitionsverweis {hasil kali silang}{}{}
dalam $\R^3$.

}
{

Kita peroleh

\mavergleichskettealign
{\vergleichskettealign
{ (N (\gamma(t)) \times F(t))'
}
{ =} { (N (\gamma(t)))' \times F(t) +  N (\gamma(t)) \times F'(t)
}
{ =} { \left(  D N   \right)_{ \gamma(t)} \left(   \gamma'(t)   \right)  \times F(t) +  N (\gamma(t)) \times F'(t)
}
{ } {
}
{ } {
}
}
{}
{}{.}
Kita harus menunjukkan bahwa vektor ini selalu tegak lurus terhadap ruang tangen $T_{\gamma(t)}Y$. Karena $F'(t)$ tegak lurus terhadap
ruang tangen, vektor tersebut bergantung linear pada vektor normal $N (\gamma(t))$, sehingga suku kanan sama dengan $0$ menurut
[[K^3/Kreuzprodukt/Eigenschaften/Fakt|Lema 33.3 (Aljabar linear (Osnabrück 2024-2025))  (3)]].
Dalam suku kiri, $F(t)$ tangensial menurut asumsi dan $\left(  D N   \right)_{ \gamma(t)} \left(   \gamma'(t)   \right)$ tangensial menurut
[[Hyperfläche/Weingartenabbildung/Tangentialraum/Fakt|Lema 4.2 (Geometri diferensial (Osnabrück 2023))]].
Karena itu, menurut
[[K^3/Kreuzprodukt/Eigenschaften/Fakt|Lema 33.3 (Aljabar linear (Osnabrück 2024-2025))  (6)]],
suku kiri tegak lurus terhadap ruang tangen, sehingga medan vektor tersebut secara keseluruhan paralel.

 }

__NOEDITSECTION__

\inputaufgabeklausurloesung
{9}

{

Misalkan

\mavergleichskette
{\vergleichskette
{ M
}
{ \neq }{ \emptyset
}
{  }{
}
{    }{
}
{   }{
}
}
{}{}{}
suatu
\definitionsverweis {manifold diferensiabel}{}{}
berdimensi $n$. Tunjukkan bahwa terdapat suatu rantai
\definitionsverweis {submanifold-submanifold tertutup}{}{}

\mavergleichskettedisp
{\vergleichskette
{ M_0
}
{ \subseteq} { M_1
}
{ \subseteq} { M_2
}
{ \subseteq  \ldots \subseteq} { M_{n-1}
}
{ \subseteq} { M_n
}
}
{
\vergleichskettefortsetzung
{ =} { M
}
{ } {}
{ } {}
{ } {}
}{}{}
sedemikian sehingga submanifold tertutup $M_i$ berdimensi $i$.

}
{

Misalkan $P \in M$ suatu titik dan
\mathl{P \in U}{} suatu lingkungan peta terbuka beserta peta
\maabbdisp {\alpha} { U } {V \subseteq \R^n
} {,}
dengan $V=B(0,3)$ bola terbuka berpusat di $0$ dan berjari-jari $3$, serta
\mathl{\alpha(P)=0}{}. Untuk
\mathl{i=1 , \ldots , n-1}{}, kita tinjau

\mathdisp {B_i = \left\{  x \in V  \mid  (x_1-1)^2+x_2^2 + \cdots + x_n^2 = 1 , \,  x_{i+2}   = 0,\, x_{i+3} = 0 , \ldots , x_n = 0       \right\}} { . }

Jadi, $B_{n-1}$ adalah permukaan bola berpusat di
\mathl{(1,0 , \ldots , 0)}{} dan berjari-jari $1$; $B_{n-2}$ di dalamnya didefinisikan oleh
\mathl{x_n=0}{} dan merupakan \anfuehrung{ekuator}{}, dan seterusnya. Kita memperoleh $B_i$ dari
\mathl{B_{i+1}}{} dengan menetapkan pula
\mathl{x_{i+2}=0}{}. Karena itu, terdapat rantai menurun subhimpunan-subhimpunan tertutup

\mathdisp {B_{n-1} \supseteq B_{n-2} \supseteq  \ldots \supseteq B_1 \supseteq B_0} {  }

\zusatzklammer {$B_0$ terdiri atas kedua titik
\mathl{(\pm1,0 , \ldots , 0 )}{}} {} {.}
Kita pandang $B_i$ sebagai serat di atas titik nol dari pemetaan
\maabbeledisp {\varphi_i} { V } {\R^{n-i}
} {(x_1 , \ldots , x_n) } {  ((x_1-1)^2+x_2^2 + \cdots + x_n^2 - 1 , x_{i+2} , \ldots , x_n)
} {,}.
Matriks Jacobi pemetaan ini adalah

\mathdisp {\begin{pmatrix}  2x_1-2 &  2x_2  &  \ldots  &  2x_{i+1} &  2x_{i+2} &  \ldots &  2x_{n-1} &  2x_n
 \\  0 &  0 &  \ldots &  0 &  1 &  0 &  \ldots &  0
 \\ \vdots  & \vdots & \vdots &  \ddots &  \ddots &  \ddots  &  \ddots  & \vdots \\  0 &  0 &  \ldots &  \ldots &  \ldots &  0 &  1 &  0 \\  0 &  0 &  \ldots &  \ldots &  \ldots  &  \ldots &  0 &  1 \end{pmatrix}} { . }

Rank matriks ini lebih kecil dari $n-i$ hanya jika
\mathl{x_1=1, \, x_2 = \cdots =x_{i+1}=0}{}, dan titik seperti itu tidak terletak pada $B_i$. Artinya, pemetaan $\varphi_i$ reguler pada serat $B_i$, sehingga berdasarkan teorema pemetaan implisit, $B_i$ merupakan submanifold tertutup dari $V$ yang berdimensi
\mathl{i}{}.

Sekarang kita tetapkan
\mathl{M_i= \alpha^{-1}(B_i)}{} untuk
\mathl{i=1 , \ldots , n-1}{} dan
\mathl{M_n= M}{.} Karena $B_i$ kompak, $M_i$ juga merupakan subhimpunan tertutup dalam $M$. Karena syarat bagi suatu manifold tertutup merupakan sifat lokal, himpunan-himpunan ini adalah submanifold-submanifold tertutup dari $M$.

 }

__NOEDITSECTION__

\inputaufgabeklausurloesung
{0}

{

}
{  }

__NOEDITSECTION__

\inputaufgabeklausurloesung
{8}

{

Misalkan
\maabbdisp {\psi} {\R^n } {\R
} {}
suatu
\definitionsverweis {fungsi diferensiabel}{}{}
dan

\mavergleichskettedisphandlinks
{\vergleichskettedisphandlinks
{M
}
{ =} { \left\{  \left(  x_1   , \,   , \ldots ,  , \,   x_n , \,   \psi(x_1 , \ldots , x_n)                \right)   \mid  (x_1 , \ldots , x_n) \in \R^n         \right\}
}
{ \subseteq} { \R^n \times \R
}
{ } {
}
{ } {
}
}
{}{}{}
merupakan
\definitionsverweis {grafik}{}{}
dari $\psi$. Tunjukkan bahwa $M$ suatu
\definitionsverweis {manifold Riemann}{}{}
yang
\definitionsverweis {berorientasi}{}{}
dan bahwa untuk
\definitionsverweis {bentuk volume kanonik}{}{}
$\omega$ pada $M$ berlaku rumus

\mavergleichskettedisp
{\vergleichskette
{ (\operatorname{Id} \times \psi)^*(\omega)
}
{ =} { \left(  1+ \sum_{i = 1}^n \left(  { \frac{   \partial   \psi  }{  \partial   x_i   } }  \right)^2  \right)^{1/2} dx_1 \wedge \ldots \wedge dx_n
}
{ } {
}
{ } {
}
{ } {
}
}
{}{}{.}

}
{

Misalkan

\mavergleichskette
{\vergleichskette
{ W
}
{ = }{ \R^n
}
{  }{
}
{    }{
}
{   }{
}
}
{}{}{.}
Pemetaan
\maabbeledisp {\varphi = \operatorname{id} \times \psi} { W } { W \times \R
} { x } {(x, \psi(x))
} {,}
merupakan suatu
\definitionsverweis {difeomorfisme}{}{}
antara $W$ dan grafik $M$. Grafik tersebut adalah
\definitionsverweis {submanifold tertutup}{}{}
dari
\mathl{W \times \R}{} dan karena itu membawa
\definitionsverweis {struktur Riemann terinduksi}{}{}
serta
\zusatzklammer {karena orientasi $W$ berpindah ke $M$} {} {}
suatu
\definitionsverweis {bentuk volume kanonik}{}{}
$\omega$. Pada situasi ini kita dapat menerapkan
[[Riemannsche Untermannigfaltigkeit des R^n/Einbettung/Volumenform/Fakt|Teorema 16.8 (Geometri diferensial (Osnabrück 2023))]].
\definitionsverweis {Turunan parsial}{}{}
dari $\varphi$ terhadap variabel ke-$i$ adalah

\mavergleichskettedisp
{\vergleichskette
{ { \frac{   \partial  \varphi  }{  \partial   x_i   } }
}
{ =} { \begin{pmatrix}  0  \\\vdots\\  0\\ 1\\  0\\ \vdots\\  0\\  { \frac{   \partial   \psi  }{  \partial   x_i   } }   \end{pmatrix}
}
{ } {
}
{ } {
}
{ } {
}
}
{}{}{.}
Misalkan

\mavergleichskette
{\vergleichskette
{Q
}
{ \in }{ W
}
{  }{
}
{    }{
}
{   }{
}
}
{}{}{}
suatu titik yang kita substitusikan ke dalam fungsi-fungsi berikut, sehingga kita menghitung dengan bilangan real di seluruh bagian. Hasil kali skalar yang membentuk entri-entri
\mathl{b_{ij}}{} dari matriks $B$
\zusatzklammer {yang determinannya harus kita hitung akarnya} {} {,}
adalah

\mavergleichskettedisp
{\vergleichskette
{ b_{ij}
}
{ =} { \left\langle  { \frac{   \partial  \varphi  }{  \partial   x_i   } }(Q)  ,  { \frac{   \partial  \varphi  }{  \partial   x_j   } }(Q)   \right\rangle
}
{ } {}
{ } {}
{ } {}
}
{}{}{}
Kita tulis

\mavergleichskette
{\vergleichskette
{B
}
{ = }{ E_n+A
}
{  }{
}
{    }{
}
{   }{
}
}
{}{}{}
dengan

\mavergleichskette
{\vergleichskette
{ A
}
{ = }{ \left(  a_{ij}  \right)_{1 \leq i,j \leq n}
}
{  }{
}
{    }{
}
{   }{
}
}
{}{}{.}
Dengan

\mavergleichskette
{\vergleichskette
{ c_i
}
{ = }{ { \frac{   \partial   \psi  }{  \partial   x_i   } } ( Q)
}
{  }{
}
{    }{
}
{   }{
}
}
{}{}{}
kita dapat menulis

\mavergleichskette
{\vergleichskette
{ a_{ij}
}
{ = }{ c_ic_j
}
{  }{
}
{    }{
}
{   }{
}
}
{}{}{}
dan keseluruhan matriks $A$ sebagai

\mavergleichskettedisp
{\vergleichskette
{A
}
{ =} { \begin{pmatrix} c_1  \\\vdots\\ c_n   \end{pmatrix} \left( c_1   , \,   \ldots  , \,  c_n                 \right)
}
{ } {
}
{ } {
}
{ } {
}
}
{}{}{}
Karena itu, $A$ merepresentasikan suatu pemetaan linear dari $\R^n$ ke $\R^n$ yang melalui $\R$
\definitionsverweis {terfaktorkan}{}{}
dan dengan demikian memiliki
\definitionsverweis {kernel}{}{,}
yang sekurang-kurangnya berdimensi $(n-1)$. Kita namakan kernel ini $K$. Jika dimensinya $n$, maka

\mavergleichskette
{\vergleichskette
{A
}
{ = }{ 0
}
{  }{
}
{    }{
}
{   }{
}
}
{}{}{}
dan $B$ adalah identitas, sehingga pernyataannya benar. Jadi, misalkan

\mavergleichskette
{\vergleichskette
{A
}
{ \neq }{ 0
}
{  }{
}
{    }{
}
{   }{
}
}
{}{}{.}
Maka

\mavergleichskette
{\vergleichskette
{v
}
{ = }{ \begin{pmatrix} c_1  \\\vdots\\ c_n   \end{pmatrix}
}
{  }{
}
{    }{
}
{   }{
}
}
{}{}{}
adalah
\definitionsverweis {vektor eigen}{}{}
dari $A$ dengan
\definitionsverweis {nilai eigen}{}{}

\mavergleichskette
{\vergleichskette
{ c_1^2 + \cdots + c_n^2
}
{ \neq }{ 0
}
{  }{
}
{    }{
}
{   }{
}
}
{}{}{.}
Vektor ini merupakan vektor eigen dari $B$ dengan nilai eigen
\mathl{1+c_1^2 + \cdots + c_n^2}{}, sedangkan $K$ membentuk
\mathl{(n-1)}{-}dimensional
\definitionsverweis {ruang eigen}{}{}
untuk $B$ dengan nilai eigen $1$. Secara keseluruhan, $B$
\definitionsverweis {dapat didiagonalkan}{}{}
dan
\definitionsverweis {determinan}{}{}
matriks tersebut merupakan hasil kali nilai-nilai eigennya, yaitu
\mathl{1+c_1^2 + \cdots + c_n^2}{.}

 }

__NOEDITSECTION__

\inputaufgabeklausurloesung
{6}

{

Hitung luas permukaan bola satuan.

}
{  }

__NOEDITSECTION__

\inputaufgabeklausurloesung
{0}

{

}
{  }

__NOEDITSECTION__

\inputaufgabeklausurloesung
{3}

{

Hitung
\definitionsverweis {turunan eksterior}{}{}
$d \omega$ dari
$1$-\definitionsverweis {bentuk diferensial}{}{}

\mavergleichskettedisp
{\vergleichskette
{ \omega
}
{ =} { { \frac{   x^2 }{  y } } dx- { \frac{   x  }{  y^2 } } dy
}
{ } {
}
{ } {
}
{ } {
}
}
{}{}{}
pada

\mavergleichskette
{\vergleichskette
{ U
}
{ = }{ \left\{ (x,y) \in \R^2  \mid   y \neq 0         \right\}
}
{  }{
}
{    }{
}
{   }{
}
}
{}{}{.}

}
{

Kita peroleh

\mavergleichskettealign
{\vergleichskettealign
{d \omega
}
{ =} { d { \frac{   x^2 }{  y } } \wedge dx - d { \frac{   x  }{  y^2 } } \wedge dy
}
{ =} { - { \frac{   x^2 }{  y^2 } } dy  \wedge dx  -   { \frac{   1 }{  y^2 } }  dx \wedge dy
}
{ =} { { \frac{   x^2 -1 }{  y^2 } }  dx \wedge dy
}
{ } {
}
}
{}
{}{.}

 }

__NOEDITSECTION__

\inputaufgabeklausurloesung
{3}

{

Berikan contoh suatu
\definitionsverweis {manifold berbatas}{}{}
yang
\definitionsverweis {terhubung}{}{,}
dan batasnya terdiri atas
\zusatzklammer {gabungan saling lepas dari} {} {}
tak hingga banyak lingkaran $S^1$.

}
{

Kita tinjau

\mavergleichskettedisp
{\vergleichskette
{ M
}
{ =} { \R^2 \setminus \biguplus_{n \in \Z} U \left(  (n,0), { \frac{   1 }{  3 } }   \right)
}
{ } {
}
{ } {
}
{ } {
}
}
{}{}{.}
Jadi, tak hingga banyak bola terbuka saling lepas dikeluarkan dari bidang real sepanjang sumbu $x$. Setiap bola menghasilkan satu komponen batas dari $M$ yang difeomorfik dengan lingkaran $S^1$. Jelas bahwa $M$ terhubung-lintasan.

 }

__NOEDITSECTION__

\inputaufgabeklausurloesung
{0}

{

}
{  }

__NOEDITSECTION__

\inputaufgabeklausurloesung
{0}

{

}
{  }

__NOEDITSECTION__

\inputaufgabeklausurloesung
{7}

{

Misalkan $M$ suatu
\definitionsverweis {manifold Riemann}{}{}
dan diberikan suatu
\definitionsverweis {koneksi linear}{}{}
pada
\definitionsverweis {bundel tangen}{}{}
$TM$ yang
\definitionsverweis {metrik}{}{}
dan
\definitionsverweis {bebas torsi}{}{.}
Tunjukkan bahwa untuk medan-medan vektor $X,Y,Z$ berlaku rumus

\mavergleichskettedisp
{\vergleichskette
{ \left\langle \nabla_X Y , Z  \right\rangle
}
{ =} { { \frac{   1  }{  2 } }  \left(  D_X \left\langle  Y  ,  Z  \right\rangle +  D_Y \left\langle  Z  ,  X  \right\rangle -  D_Z \left\langle  X  ,  Y  \right\rangle  - \left\langle  Y  , [X,Z]   \right\rangle - \left\langle  Z  , [Y,X]   \right\rangle + \left\langle  X  , [Z,Y]   \right\rangle   \right)
}
{ } {
}
{ } {
}
{ } {
}
}
{}{}{.}

}
{

Karena koneksi bebas torsi, berlaku

\mavergleichskettedisp
{\vergleichskette
{ \nabla_VW- \nabla_WV
}
{ =} { [V,W]
}
{ } {
}
{ } {
}
{ } {
}
}
{}{}{,}

\mavergleichskettedisp
{\vergleichskette
{ \nabla_VZ- \nabla_ZV
}
{ =} { [V,Z]
}
{ } {
}
{ } {
}
{ } {
}
}
{}{}{}
dan

\mavergleichskettedisp
{\vergleichskette
{ \nabla_WZ- \nabla_ZW
}
{ =} { [W,Z]
}
{ } {
}
{ } {
}
{ } {
}
}
{}{}{,}
dengan identitas pertama juga dapat ditulis sebagai

\mavergleichskettedisp
{\vergleichskette
{ \nabla_VW + \nabla_WV
}
{ =} { 2 \nabla_VW -[V,W]
}
{ } {
}
{ } {
}
{ } {
}
}
{}{}{}.
Karena koneksi metrik, berlaku identitas-identitas

\mavergleichskettedisp
{\vergleichskette
{ D_V \left\langle  W  ,  Z  \right\rangle
}
{ =} { \left\langle \nabla_V W , Z  \right\rangle  + \left\langle  W  ,  \nabla_V Z  \right\rangle
}
{ } {
}
{ } {
}
{ } {
}
}
{}{}{,}

\mavergleichskettedisp
{\vergleichskette
{ D_W \left\langle  Z  ,  V  \right\rangle
}
{ =} { \left\langle \nabla_WZ , V  \right\rangle  + \left\langle  Z  ,  \nabla_WV  \right\rangle
}
{ } {
}
{ } {
}
{ } {
}
}
{}{}{}
dan

\mavergleichskettedisp
{\vergleichskette
{ D_Z \left\langle  V  ,  W  \right\rangle
}
{ =} { \left\langle \nabla_Z V , W  \right\rangle  + \left\langle  V  ,  \nabla_Z W   \right\rangle
}
{ } {
}
{ } {
}
{ } {
}
}
{}{}{.}
Kita jumlahkan persamaan-persamaan tersebut dan memperoleh

\mavergleichskettedisphandlinks
{\vergleichskettedisphandlinks
{ D_V \left\langle  W  ,  Z  \right\rangle+ D_W \left\langle  Z  ,  V  \right\rangle  - D_Z \left\langle  V  ,  W  \right\rangle
}
{ =} { \left\langle \nabla_V W+\nabla_W V  ,  Z  \right\rangle + \left\langle \nabla_W Z-\nabla_Z W   ,  V  \right\rangle+\left\langle  \nabla_V Z-\nabla_Z V   ,  W  \right\rangle
}
{ } {
}
{ } {
}
{ } {
}
}
{}{}{.}
Kita substitusikan bentuk-bentuk yang diperoleh di atas ke ruas kanan dan mendapatkan

\mavergleichskettealignhandlinks
{\vergleichskettealignhandlinks
{ D_V \left\langle  W  ,  Z  \right\rangle+ D_W \left\langle  Z  ,  V  \right\rangle  - D_Z \left\langle  V  ,  W  \right\rangle
}
{ =} { \left\langle  2 \nabla_VW -[V,W]  ,  Z  \right\rangle + \left\langle  [W,Z]  ,  V  \right\rangle + \left\langle [V,Z]   ,  W  \right\rangle
}
{ =} { 2\left\langle  \nabla_VW   ,  Z  \right\rangle - \left\langle  [V,W]  ,  Z  \right\rangle + \left\langle  [W,Z]  ,  V  \right\rangle + \left\langle  [V,Z]   ,  W  \right\rangle
}
{ } {
}
{ } {
}
}
{}
{}{.}
Dengan menata ulang persamaan ini, diperoleh rumus yang diminta.

Berdasarkan persamaan tersebut,
\mathl{\left\langle  \nabla_V W , Z  \right\rangle}{} untuk sebarang medan vektor ditentukan oleh ruas kanan, yang sama sekali tidak memuat koneksi. Karena ini berlaku untuk setiap medan vektor $Z$, turunan vertikal $\nabla_V W$ dan dengan demikian koneksi linearnya juga ditentukan secara unik.

 }

 [[Kategorie:Latexseite]]
